The cryptocurrency market has grown from a single asset—Bitcoin, introduced by
\citet{nakamoto2008bitcoin} as a decentralized peer-to-peer payment system—into
a diverse asset class encompassing thousands of tokens with heterogeneous
on-chain economics, governance structures, and macro exposures.
\citet{liu2022cryptocurrency} document that the cross-section of cryptocurrency
returns exhibits its own set of risk factors, including momentum and size premia
that are structurally distinct from their equity counterparts.
With over 86 actively traded tokens spanning divergent network fundamentals
and social sentiment signals, predicting which assets will outperform their peers
in the following week is a natural cross-sectional learning problem—one that
demands models capable of synthesizing heterogeneous feature sets while learning
a \emph{relative}, rather than absolute, notion of return.

\paragraph{Prior work and the research gap.}
The foundational study of \citet{gu2020empirical} demonstrates that machine
learning models—in particular, tree ensembles and deep networks—substantially
outperform linear baselines for cross-sectional equity return prediction when
feature sets are rich and model capacity is appropriately regularized.
\citet{freyberger2020dissecting} extend this finding to a nonparametric
setting, showing that nonlinear interactions among stock characteristics carry
significant incremental predictive power.
Applying these insights to cryptocurrency markets is appealing in principle
but nontrivial in practice.
The primary obstacle is \emph{cross-sectional volatility heterogeneity}: a
Bitcoin weekly return might span $\pm 5\%$ while a small-cap altcoin regularly
moves $\pm 30\%$.
When the standard ListNet loss \citep{cao2007learning} is trained with raw
weekly returns as softmax targets, the resulting target distribution $Q$ is
dominated in volatile weeks by a handful of extreme-return assets, and nearly
uniform in calm weeks.
This instability produces high-variance gradient estimates that undermine
consistent ranking performance across market regimes.
\citet{poh2021building} identify this failure mode in equity cross-sectional
strategies and propose replacing raw-return targets with cross-sectional rank
percentiles; we adopt and extend their framework to the cryptocurrency setting,
where the problem is more acute due to the larger cross-sectional dispersion
in volatility.

On the data side, \citet{liu2022cryptocurrency} and \citet{cong2021tokenomics}
establish that on-chain network metrics—such as network value-to-transaction
ratio (NVT) and market-value-to-realized-value (MVRV)—contain information
about the fundamental valuation state of individual tokens that is orthogonal
to price-based momentum signals.
Beyond on-chain data, the role of alternative information sources in asset
pricing has been well established in the equity literature.
\citet{tetlock2007giving} demonstrates that media pessimism predicts short-run
downward price pressure in equity markets, and \citet{bollen2011twitter}
show that Twitter-derived mood measures Granger-cause changes in the Dow Jones
Industrial Average.
We extend this line of inquiry to a novel source: the social media activity of
a single, highly influential policy actor (the 45th and 47th U.S.\ President),
whose public communications on cryptocurrency regulation and macroeconomic
tariff policy have preceded market-wide repricing events.
A further structural data source—unavailable before January 2024—is the
granular daily flow data for spot Bitcoin and Ethereum exchange-traded funds
published by Farside Investors
\citep{bitcoin_etf_farside,ethereum_etf_farside}, which permits decomposition of ETF
inflows into net new capital and legacy-holder rotations.
To the best of our knowledge, this paper is the first to integrate on-chain fundamentals,
policy-linked social media signals, and ETF structural flows into a unified
cross-sectional ranking framework evaluated on a live 86-token universe.

\paragraph{Architecture considerations.}
Two families of temporal deep learning architectures are well suited to
sequential financial data.
The Long Short-Term Memory network \citep{hochreiter1997long} encodes the
$L$-step feature history into a fixed-size hidden state via gated recurrence,
and has remained a competitive baseline for financial time series despite the
rise of attention-based models.
The Temporal Fusion Transformer \citep{lim2021temporal} combines interpretable
Variable Selection Networks (VSN) with multi-head self-attention
\citep{vaswani2017attention}, enabling explicit feature selection at each
time step and scalable handling of heterogeneous input types.
However, \citet{gu2020empirical} also observe that cross-sectional features
can dominate temporal features for equity ranking—motivating a purely
cross-sectional architecture that operates on contemporaneous inputs alone,
without any temporal encoder.
We formalize and test this design hypothesis through our CrossSectionalGatedNet.

\paragraph{Contributions.}
Building on the above foundation, this paper makes four targeted contributions
to the cross-sectional cryptocurrency prediction literature:
\begin{enumerate}
  \item \textbf{Rank-normalized ListNet targets.}
        Following \citet{poh2021building}, we replace raw returns with
        cross-sectional rank percentiles $\in [0,1]$ as softmax targets,
        ensuring equal contribution of each asset regardless of its
        absolute volatility and stabilizing the cross-entropy gradient
        across heterogeneous market regimes.
  \item \textbf{Gradient accumulation.}
        The effective cross-sectional batch size is increased from 86
        (one week's cross-section) to $86 \times 4 = 344$ assets per
        gradient step, providing a more stable rank estimate for the
        ListNet loss and reducing stochasticity in the gradient direction.
  \item \textbf{CrossSectionalGatedNet.}
        A new architecture consisting purely of Gated Residual Networks
        (GRN) and cross-asset attention, with no temporal encoder.
        Motivated by \citet{gu2020empirical}, we test whether the
        cross-sectional structure alone—without temporal modeling—is
        sufficient for the ranking task when on-chain fundamentals
        dominate the predictive signal.
  \item \textbf{Full-scale TFT Variable Selection Network.}
        We implement a per-variable shared GRN ($\sim$27K parameters) for
        nonlinear variable selection, closer to the original
        \citet{lim2021temporal} specification, and show that this design
        substantially stabilizes feature selection.
\end{enumerate}

\paragraph{Paper organization.}
\Cref{sec:related} reviews related work in greater detail.
\Cref{sec:data} describes the 86-token, 324-week balanced panel and its
seven feature groups.
\Cref{sec:method} formalizes the ranking loss and the four model architectures.
\Cref{sec:experiments} details the training setup, six feature configurations,
and 32-seed ensemble strategy.
\Cref{sec:results} presents the main cross-model comparison.
\Cref{sec:analysis} provides ablation studies, Granger causality tests,
transaction cost sensitivity analysis, and sub-period robustness checks.
\Cref{sec:conclusion} concludes.
